\chapter{Experimental Setup and System Configuration}
All reported energy runs used NVIDIA RTX 6000 Ada GPUs on one shared SLURM node at Sardar Patel Institute of Technology, with the same 500 prompts, greedy decoding and seed 42. In this chapter, ``not documented'' means that the repository records the value but not the reason for it, and ``not recorded'' means that the value itself is absent.

\section{Hardware}

\begingroup\footnotesize\setlength{\tabcolsep}{4pt}
\begin{longtable}{>{\raggedright\arraybackslash}p{25.2mm}>{\raggedright\arraybackslash}p{39.8mm}>{\raggedright\arraybackslash}p{43.7mm}>{\raggedright\arraybackslash}p{30.1mm}}
\toprule
\textbf{Component} & \textbf{Specification} & \textbf{Reason and contribution} & \textbf{Evidence} \\
\midrule
\endhead
GPU & NVIDIA RTX 6000 Ada Generation, 49,140 MiB, 300 W cap & The college cluster replaced Kaggle after a Kaggle session hung until its 12-hour limit. The GPU exposes the NVML energy counter, and vLLM lists FP8 support on Ada~\cite{b45}. & nvidia-smi logs; NEW.md \\
GPU units & Two on the node: GPU-c53791fb (bus 01:00) and GPU-da9c4cd1 (bus 61:00) & Assigned by the scheduler, not chosen. Recorded because the unit is confounded with time of day (section 5.5). & Run info files \\
CPU & 2 x AMD EPYC 9554, reported as 56 cores per socket, 2 threads per core, 224 logical CPUs & Part of the shared node. Its load is relevant to the host-side slowdown. & lm-eval environment record \\
Memory & 250 GB RAM & Part of the shared node. & CLUSTER\_\allowbreak{}MANUAL.md \\
Job allocation & 1 GPU, 8 CPUs; 32 GB (Stack A), 64 GB (Stack B) & Values not documented, beyond limiting thread pools to the granted CPUs. & sbatch scripts \\
Adapter training & Kaggle, Tesla T4 & The platform available when the adapters were trained. & Notebook output \\
\bottomrule
\end{longtable}
\endgroup

\section{Software}

\begingroup\footnotesize\setlength{\tabcolsep}{4pt}
\begin{longtable}{>{\raggedright\arraybackslash}p{33.3mm}>{\raggedright\arraybackslash}p{33.1mm}>{\raggedright\arraybackslash}p{30.7mm}>{\raggedright\arraybackslash}p{41.7mm}}
\toprule
\textbf{Software} & \textbf{Stack A} & \textbf{Stack B} & \textbf{Reason and contribution} \\
\midrule
\endhead
Operating system & Rocky Linux 9.8 (same node) & Rocky Linux 9.8 & Cluster default \\
NVIDIA driver & 610.43.02 & 610.57.04 & Cluster-managed; the change between stacks is a stated limitation \\
PyTorch, CUDA & 2.13.0+cu130, CUDA 13.0 & 2.13.0+cu130, CUDA 13.0 & Execution framework for both stacks \\
Python & 3.11 & 3.11.14 & Virtual environments built with uv \\
Transformers & Not recorded & 5.17.0 & Stack A uses it as the common drop-in inference path \\
bitsandbytes & Not recorded & Not used for the tiers & Provides the NF4 and LLM.int8() paths under test \\
vLLM & Not used & 0.31.0, CUDA graphs on & Native FP8 and Marlin AWQ kernels, used to test whether the Stack A result depends on implementation \\
lm-eval-harness & 0.4.13 & 0.4.13 & Community-standard, versioned benchmarks \\
PEFT & 0.19.1, recorded by the adapters & Not used & Loads LoRA adapters \\
NVML bindings & pynvml, version not recorded & nvidia-ml-py, version not recorded & Reads the energy counter directly \\
Router libraries & scikit-fuzzy, textstat, spaCy en\_\allowbreak{}core\_\allowbreak{}web\_\allowbreak{}sm; versions not recorded & Same modules, replayed offline & Mamdani inference, readability grade and dependency depth \\
Analysis libraries & pandas, NumPy, SciPy, scikit-learn, Matplotlib; cluster versions not recorded & Same & Bootstrap, replay, correlation and figures \\
\bottomrule
\end{longtable}
\endgroup

Versions marked ``not recorded'' were not captured at the time of the runs. The file requirements.txt predates these runs and does not describe them.

\section{Models and configurations}

\begingroup\footnotesize\setlength{\tabcolsep}{4pt}
\begin{longtable}{>{\raggedright\arraybackslash}p{27.5mm}>{\raggedright\arraybackslash}p{33.4mm}>{\raggedright\arraybackslash}p{27.1mm}>{\raggedright\arraybackslash}p{17.3mm}>{\raggedright\arraybackslash}p{30.7mm}}
\toprule
\textbf{Configuration} & \textbf{Model} & \textbf{Quantization} & \textbf{Adapter} & \textbf{Reason} \\
\midrule
\endhead
Stack A float16 & Llama-3.2-1B~\cite{b36} & None & r=4 attached & Reference tier. The model matches the adapters and fit the earlier T4 target. \\
Stack A 8-bit & Llama-3.2-1B & LLM.int8() & r=8 attached & Common drop-in 8-bit path; computes in 16-bit~\cite{b6} \\
Stack A 4-bit & Llama-3.2-1B & NF4, double quantization, float16 compute & r=16 attached & Common drop-in 4-bit path. Double quantization: not documented. \\
Stack B 1B float16 & Llama-3.2-1B & None & None & Continuity with Stack A \\
Stack B 1B FP8 & Llama-3.2-1B & fp8 & None & Native FP8 on the same model \\
Stack B 8B float16 & Llama-3.1-8B-Instruct~\cite{b36} & None & None & A larger model, expected to be more memory-bound in decoding \\
Stack B 8B FP8 & Llama-3.1-8B-Instruct & fp8 & None & Same comparison at 8B \\
Stack B 8B AWQ & hugging-quants/Meta-Llama-3.1-8B-Instruct-AWQ-INT4 & AWQ, resolved as auto\_\allowbreak{}awq & None & Public 4-bit checkpoint with Marlin kernels \\
\bottomrule
\end{longtable}
\endgroup

The 1B model is a base model and the 8B model is instruction-tuned, so 1B versus 8B accuracy is not a pure comparison of size. No 1B AWQ configuration was run.

\section{Experimental parameters}

\begingroup\footnotesize\setlength{\tabcolsep}{4pt}
\begin{longtable}{>{\raggedright\arraybackslash}p{32.8mm}>{\raggedright\arraybackslash}p{50.3mm}>{\raggedright\arraybackslash}p{58.6mm}}
\toprule
\textbf{Parameter} & \textbf{Value} & \textbf{Reason and contribution} \\
\midrule
\endhead
Decoding & Greedy (do\_\allowbreak{}sample false; temperature 0) & Determinism allows one grid to serve all policies \\
Maximum new tokens & 128 & Fixed cap across configurations; the value is not documented \\
Warm-up & 5 discarded generations & Excludes start-up effects; the count is not documented \\
Batch size & 1 (Stack A); 1 and 8, static batches of consecutive prompts (Stack B) & Batch 1 represents single-request decoding; batch 8 tests whether savings persist \\
Minimum tokens for Stack A J/token & 16 & Excludes generations too short for the counter to advance \\
Seed & 42 & Reproducibility \\
Stack B engine & dtype float16, gpu\_\allowbreak{}memory\_\allowbreak{}utilization 0.85, max\_\allowbreak{}model\_\allowbreak{}len 1024 (energy), 2048 (accuracy) & Values not documented \\
Monitor interval & 2 s & Continuous sampling replaces the failed end-of-run snapshot; the interval is not documented \\
Watchdog & 900 s with no progress & Added after one configuration hung while loading \\
Disruption flag & Any co-tenant, or throughput over 20\% below the configuration's median across runs & Targets the observed signature; the 20\% value is a judgement, not derived \\
Stack A thread pools & OMP, MKL, OpenBLAS and NumExpr limited to 8 & Prevents oversubscription against the granted CPUs \\
Tier cuts & 33 and 66 on a 0 to 100 score & Equal division of the scale; tested in section 5.6 \\
Bootstrap & 2,000 resamples & Interval estimation without a normality assumption \\
\bottomrule
\end{longtable}
\endgroup


\begingroup\footnotesize\setlength{\tabcolsep}{4pt}
\begin{longtable}{>{\raggedright\arraybackslash}p{44.7mm}>{\raggedright\arraybackslash}p{53.0mm}>{\raggedright\arraybackslash}p{43.9mm}}
\toprule
\textbf{LoRA hyperparameter} & \textbf{Value} & \textbf{Reason} \\
\midrule
\endhead
Rank, alpha & 16/32 (4-bit tier), 8/16 (8-bit tier), 4/8 (16-bit tier) & Lower precision receives more capacity; constant alpha/r of 2 \\
Target modules & q\_\allowbreak{}proj, v\_\allowbreak{}proj & Not documented \\
Dropout & 0.05 & Not documented \\
Training data & Up to 2,000 samples per complexity class, maximum length 256 & Not documented \\
Training & 1 epoch, batch size 4, learning rate 2e-4, paged\_\allowbreak{}adamw\_\allowbreak{}8bit, seed 42 & Not documented \\
\bottomrule
\end{longtable}
\endgroup

\section{Datasets}
The evaluation set has 500 unique prompts, built by prepare\_\allowbreak{}eval\_\allowbreak{}dataset.py with seed 42. The builder fails rather than pads with duplicates, because an earlier set contained only 30 unique prompts among 50 rows. Difficulty labels record the source dataset and do not measure difficulty for a given model. The 200/150/150 split is not otherwise justified in the repository.


\begingroup\footnotesize\setlength{\tabcolsep}{4pt}
\begin{longtable}{>{\raggedright\arraybackslash}p{16.8mm}>{\raggedright\arraybackslash}p{17.7mm}>{\raggedright\arraybackslash}p{41.5mm}>{\raggedright\arraybackslash}p{28.1mm}>{\raggedright\arraybackslash}p{31.8mm}}
\toprule
\textbf{Label} & \textbf{Prompts} & \textbf{Source} & \textbf{Reference answer} & \textbf{Reason} \\
\midrule
\endhead
Easy & 200 & TriviaQA, rc.nocontext, validation~\cite{b38} & Short factual answer & Unambiguous answers \\
Medium & 150 & Alpaca, instructions with no input~\cite{b42} & Reference output & Explanation and description tasks \\
Hard & 75 & GSM8K, test split~\cite{b39} & Final numeric answer & Multi-step arithmetic \\
Hard & 75 & CodeAlpaca-20k~\cite{b46} & Reference output & Coding tasks \\
\bottomrule
\end{longtable}
\endgroup

Benchmark accuracy uses tinyBenchmarks~\cite{b37}, with 100 curated items each from MMLU~\cite{b40}, HellaSwag~\cite{b41} and GSM8K~\cite{b39}, selected to reduce GPU time. Scores are IRT-estimated, and lm-eval reports no standard error for them. The adapter training data, combined\_\allowbreak{}cleaned.csv, is not included in the repository.

\section{Experimental runs}
Times are IST start times.


\begingroup\footnotesize\setlength{\tabcolsep}{4pt}
\begin{longtable}{>{\raggedright\arraybackslash}p{28.6mm}>{\raggedright\arraybackslash}p{14.7mm}>{\raggedright\arraybackslash}p{19.8mm}>{\raggedright\arraybackslash}p{22.1mm}>{\raggedright\arraybackslash}p{19.2mm}>{\raggedright\arraybackslash}p{28.7mm}}
\toprule
\textbf{Session} & \textbf{Job} & \textbf{Date and time} & \textbf{GPU unit} & \textbf{Adapters} & \textbf{Use in this thesis} \\
\midrule
\endhead
1, energy & Not in log & 4 Sep 2026 & Not recorded & No & Pitfall evidence; 16-bit value withdrawn \\
4, run 1 & 1505 & 5 Sep & c53791fb & Yes & Before router fix; contended \\
4, run 2, interleaved & 1507 & 5 Sep & c53791fb & Yes & Before router fix; contended \\
4, run 3 & 1529 & 9 Sep & c53791fb & Yes & Contended; not used for results \\
4, run 4 & 1691 & 29 Sep, 11:48 & c53791fb & Yes & 16-bit phase disrupted \\
4, run 5 & 1692 & 30 Sep, 00:05 & da9c4cd1 & Yes & Clean \\
Diagnostic, day & 1706 & 30 Sep, 13:50 & Not recorded & No & Host-side evidence \\
4, run 6 & 1707 & 1 Oct, 00:58 & da9c4cd1 & Yes & Clean \\
Diagnostic, night & 1732 & 3 Oct, 00:55 & Not recorded & No & Host-side evidence \\
4, run 7 & 1733 & 3 Oct, 00:57 & da9c4cd1 & Yes & Clean \\
2, Stack A accuracy & 1784 & 7 Oct & Not read & With and without & Accuracy per tier \\
5, 1B and AWQ accuracy & 2062 & 27 Oct & Not read & No & Stack B accuracy \\
5, 8B accuracy & 2067 & 27 Oct & Not read & No & Stack B accuracy \\
5, run 1 & 2066, 2093 & 28 Oct, 11:02; AWQ batch 8 rerun 29 Oct & c53791fb & No & Clean \\
5, run 2 & 2094 & 29 Oct, 14:23 & c53791fb & No & Clean \\
5, run 3 & 2095 & 29 Oct, 15:19 & c53791fb & No & Clean \\
\bottomrule
\end{longtable}
\endgroup

\section{Evaluation metrics}

\begingroup\footnotesize\setlength{\tabcolsep}{4pt}
\begin{longtable}{>{\raggedright\arraybackslash}p{50.8mm}>{\raggedright\arraybackslash}p{26.2mm}>{\raggedright\arraybackslash}p{64.6mm}}
\toprule
\textbf{Metric} & \textbf{Unit} & \textbf{Purpose} \\
\midrule
\endhead
Energy per token & J/token & Main efficiency measure, defined per stack in section 3.4 \\
Energy per request & J/request & Cost of serving a prompt, used for routing \\
Energy per correct answer & J & Combines cost and proxy accuracy; descriptive only \\
Throughput & tokens/s & Explains energy and detects disruption \\
Mean power & W & Separates working from waiting \\
Benchmark accuracy & 0 to 1 & tinyMMLU, tinyHellaswag acc\_\allowbreak{}norm; tinyGSM8k flexible and strict \\
Proxy correctness & 0 to 1 & Reference match on the 500 prompts, for paired comparison only \\
Same-answer share & \% & Share of prompts answered identically to float16 \\
Cross-validated AUC & 0 to 1 & Whether the features predict a target; 0.5 is chance \\
\bottomrule
\end{longtable}
\endgroup

\section{Testing conditions and deviations from plan}
Stack A checked co-tenants with one end-of-run snapshot, and Stack B sampled them continuously. All clean Stack A runs were overnight, and all Stack B runs were in daytime. Energy is GPU-board energy only.


\begingroup\footnotesize\setlength{\tabcolsep}{4pt}
\begin{longtable}{>{\raggedright\arraybackslash}p{49.3mm}>{\raggedright\arraybackslash}p{40.2mm}>{\raggedright\arraybackslash}p{52.1mm}}
\toprule
\textbf{Plan} & \textbf{What happened} & \textbf{Consequence} \\
\midrule
\endhead
Kaggle T4 for all sessions & Moved to the college cluster after a Kaggle hang & Earlier T4 assumptions do not apply \\
Stack B on three nights & All three runs were in daytime on 28-29 October; the reason is not recorded & No night comparison; confound remains \\
Feature ablation study & Not run & Feature importance is inferred from correlations only \\
Tier-preference router trained on Session 4 data & Not built & A simple logistic baseline was used instead \\
FrugalGPT-style cascade & Optional; not built beyond a stub & Excluded from results \\
APGR and CPT routing metrics & Not reported & Random-matched and oracle comparisons used instead \\
\bottomrule
\end{longtable}
\endgroup

\section{Reproducibility}
From the major-project directory, on a cluster with the same SLURM partition:

\begin{enumerate}
\item Build the prompt set: python training/scripts/prepare\_\allowbreak{}eval\_\allowbreak{}dataset.py --num-prompts 500 --seed 42 --output backend/src/green\_\allowbreak{}weight/data/eval\_\allowbreak{}prompts.jsonl.
\item Run Stack A on separate days: sbatch --export=ALL,RUN\_\allowbreak{}ID=\textless{}n\textgreater{} training/scripts/session4\_\allowbreak{}full\_\allowbreak{}run.sh. Copy each run's output files to routing\_\allowbreak{}run\textless{}n\textgreater{}\_\allowbreak{}* in major-project, because the analysis expects those names.
\item Run Stack A accuracy: sbatch training/scripts/session2\_\allowbreak{}accuracy.sh.
\item Set up Stack B with bash training/scripts/setup\_\allowbreak{}vllm\_\allowbreak{}env.sh, pinning vLLM to 0.31.0 because the script installs the latest version. Then run sbatch --export=ALL,RUN\_\allowbreak{}ID=\textless{}n\textgreater{} training/scripts/session5\_\allowbreak{}vllm.sh.
\item Run Stack B accuracy: sbatch training/scripts/session5\_\allowbreak{}accuracy.sh.
\item Analyse Stack A with python training/scripts/paper\_\allowbreak{}a\_\allowbreak{}analysis.py --clean-runs 5 6 7. Analyse Stack B by passing the run directories to session5\_\allowbreak{}analysis.py, and to session5\_\allowbreak{}router\_\allowbreak{}replay.py with --pool 1b\_\allowbreak{}fp16,8b\_\allowbreak{}awq4 --pool 8b\_\allowbreak{}awq4,8b\_\allowbreak{}fp8,8b\_\allowbreak{}fp16 --batch-size 1.
\item Run python training/scripts/thesis\_\allowbreak{}extra\_\allowbreak{}analysis.py.
\end{enumerate}
The scripts source \textasciitilde{}/greenweight\_\allowbreak{}env.sh and \textasciitilde{}/vllm\_\allowbreak{}env.sh, which are not in the repository, and the Llama models are gated. A Hugging Face token must be supplied through the environment and must not be stored in the repository.

